% !TeX root = main.tex
% ============================================================
%  摘要页  ——  评委的第一印象，权重远超它的篇幅
% ============================================================
%  通知四(一)4 明文要求，摘要必须包含：
%     建模思路 / 主要方法 / 模型 / 结果与结论 / 创新点
%  篇幅 ≤ 2 页，无需译成英文。
%
%  写作要点：
%   · 一定要有具体数字。
%     「求得最优方案总成本 3.27 万元，较基准方案下降 18.4%」
%     远胜「得到了较好的结果」。
%   · 一问一段，结构与正文对应。
%   · 最后写，写三遍。
% ============================================================

\begin{abstract}

% ---- 开篇：背景 + 本文做了什么（2~3 句） ----
针对……问题，本文……。

% ---- 问题一 ----
\textbf{针对问题一}，建立了……模型。首先……，其次……，采用……求解，
得到……。结果表明……（\textbf{给出具体数字}）。

% ---- 问题二 ----
\textbf{针对问题二}，……。

% ---- 问题三 ----
\textbf{针对问题三}，……。

% ---- 检验与灵敏度 ----
最后对模型进行了……检验与灵敏度分析，结果显示……，说明模型具有较好的稳健性。

% ---- 创新点（务必单列，评委会找） ----
\textbf{本文的创新点在于}：（1）……；（2）……；（3）……。

\keywords{关键词一\quad 关键词二\quad 关键词三\quad 关键词四}

\end{abstract}
